% Generated draft: requires mathematical and source review.
\documentclass[11pt]{article}
\usepackage[margin=25mm]{geometry}
\usepackage{fontspec}
\setmainfont{Latin Modern Roman}
\usepackage{amsmath,amssymb,mathtools,unicode-math}
\setmathfont{Latin Modern Math}
\usepackage{enumitem,needspace,xurl,hyperref}
\setlength{\parindent}{0pt}
\setlength{\parskip}{5pt}
\setlength{\emergencystretch}{3em}
\newcommand{\N}{\mathbb{N}}
\newcommand{\Z}{\mathbb{Z}}
\newcommand{\Q}{\mathbb{Q}}
\newcommand{\R}{\mathbb{R}}
\newcommand{\C}{\mathbb{C}}
\newcommand{\F}{\mathbb{F}}
\newcommand{\A}{\mathbb{A}}
\newcommand{\PP}{\mathbb{P}}
\newcommand{\E}{\mathbb{E}}
\newcommand{\eps}{\varepsilon}
\newcommand{\dd}{\,\mathrm{d}}
\begin{document}
\section*{UnsolvedMath ID: 800333}
\textbf{Source problem code:} catalogue entry
\subsection{Definitions and mathematical statement}
A unit triangle is an elementary triangular face of the triangular lattice generated by \(u=(1,0)\) and \(v=(1/2,\sqrt{3}/2)\). A lozenge is the union of two unit triangles sharing a full edge. A lozenge tiling of a finite collection of unit triangles is a partition of that collection into lozenges.

The supplied statement asks for the number of lozenge tilings of a regular hexagon after removing two of its ``innermost'' unit triangles, one upward-pointing and one downward-pointing. It does not explicitly define ``innermost'' or specify the relative positions of the removed triangles.

For a precise conditional formulation, let \(n\in\Z_{\geq 1}\), and let \(H_n\) be the regular lattice hexagon with vertices
\[
nu,\quad nv,\quad n(v-u),\quad -nu,\quad -nv,\quad n(u-v).
\]
Let \(\mathcal{T}_n\) be the collection of unit triangles contained in \(H_n\). Interpret the innermost triangles as the six unit triangles having the origin as a vertex. Choose an upward-pointing triangle \(T_+\) and a downward-pointing triangle \(T_-\) among these six. Define
\[
M_n(T_+,T_-)
=\#\bigl\{\text{lozenge tilings of }\mathcal{T}_n\setminus\{T_+,T_-\}\bigr\}.
\]
The corresponding enumeration problem is to determine \(M_n(T_+,T_-)\) as a function of the side length and the prescribed removed pair. This is a conditional interpretation of the source: the supplied excerpt does not establish that these six triangles are intended, nor does it specify whether the removed triangles share an edge or lie opposite each other around the center. Those choices must remain unresolved in identifying the intended enumeration problem.
\subsection{Short English statement}
Propp asks how many lozenge tilings a regular hexagon has when two innermost unit triangles of opposite orientations are removed. The formulation above gives a precise conditional counting problem, but the supplied excerpt leaves the meaning of ``innermost'' and the relative positions of the holes unspecified. The registry labels the problem open; that label is supplied metadata, not an independently verified assessment.
\subsection{Sources}
\begin{itemize}
\item Supplied local registry record: permanent UnsolvedMath ID \(800333\), source problem code \texttt{catalogue entry}, title ``Enumeration of Matchings --- Problem 4.'' Attribution: Propp; proposed year: 1998. The supplied registry status is \texttt{open}.
\item Propp, \emph{Twenty Open Problems in Enumeration of Matchings} (1998), Problem 4 (\#4). Supplied source URL: \url{https://arxiv.org/abs/math/9801060}. The record reports access on 2026-10-16 and extraction by \texttt{arxiv\_latex}; these are supplied provenance details, not a claim of independent source consultation.
\end{itemize}
\end{document}
